\documentclass{article}
\usepackage{fuzz}
\makeatletter
\newenvironment{theorem}[1]{\medskip\noindent\textbf{Theorem (#1).}\@nameuse{zed}}{\@nameuse{endzed}}
\makeatother

\begin{document}

\title{A Formal Specification of \texttt{model-provider-fleet} (the OpenAI-protocol harnesses: codex, pi, opencode)}
\author{Punt Labs}
\maketitle

\section{Purpose and scope}

This is one of a two-tool split. \texttt{model-provider-fleet} configures the
three AI coding harnesses that speak the OpenAI protocol --- codex, pi, opencode
--- from one operator selection. Claude Code is handled by the separate
\texttt{model-provider-claude} tool and is deliberately absent here; everything
that made the combined tool hard (canonicalization, the override map, tier-alias
fan-out, the subscription-credential path) belonged to Claude Code and is gone
from this one.

The resulting guarantee is uniform and strong: when routed through a gateway,
\emph{every} request every harness issues --- main loop and every side-request
--- puts the operator's selection on the wire, verbatim, reaching the provider
the selection names. There is no per-request divergence to reason about: no
harness here canonicalizes, so the wire id is always the selection itself.
Preconditions are explicit so TTF partition analysis can derive test cases from
each operation's input domain.

\section{Given sets and enumerations}

A model id is an opaque catalog token (e.g. \texttt{openrouter/openai/gpt-6}). A
provider is the upstream a gateway routes a model to.

\begin{zed}
  [ModelId, Provider]
\end{zed}

The three OpenAI-protocol harnesses.

\begin{zed}
  Harness ::= codex | pi | opencode
\end{zed}

Request kinds. \emph{main} is the primary loop; \emph{side} is a background /
subagent / small-model request (codex and pi sub-agents, opencode's small/title
model). Side-requests are first-class: they are the paths that fail silently
when only the main model is set.

\begin{zed}
  ReqKind ::= main | side
\end{zed}

The gateways in use; the specification is written over an abstract gateway and
these are the concrete instances. $direct$ is modelled separately (the absence
of a gateway), not as a member.

\begin{zed}
  Gateway ::= bifrost | envoy | agentgateway
\end{zed}

Booleans. fuzz's toolkit has no built-in boolean; $yes$/$no$ are its values.

\begin{zed}
  Bool ::= yes | no
\end{zed}

\section{Harness properties}

Every harness here speaks the OpenAI protocol and issues a main request and a
side request. (A harness with no distinct side model still has the slot; it
simply inherits the main model, which on the wire is the selection.)

\begin{axdef}
  issues : Harness \fun \power ReqKind
\where
  issues = \{ codex \mapsto \{ main, side \}, \\
  \t1         pi \mapsto \{ main, side \}, \\
  \t1         opencode \mapsto \{ main, side \} \}
\end{axdef}

Authentication here is the per-host gateway virtual key only. There is no
subscription-credential mode: a Claude Pro/Max subscription is licensed for
Anthropic's own clients, so it has no legitimate home on any harness in this
tool. Modelling the auth mode as a one-value type records that simplification.

\begin{zed}
  AuthMode ::= vkauth
\end{zed}

\section{The gateway contract}

For each gateway: $gwCatalog$ gives, per model id, the provider it routes to;
$gwServes$ is the subset actually entitled for this account; $gwOpenai$ records
whether the gateway exposes the OpenAI protocol these harnesses speak. Every
gateway entitles at least one id.

\begin{axdef}
  gwCatalog : Gateway \fun (ModelId \pfun Provider) \\
  gwServes : Gateway \fun \power ModelId \\
  gwOpenai : Gateway \fun Bool
\where
  \forall g : Gateway @ \\
  \t1 gwServes~g \subseteq \dom (gwCatalog~g) \land \\
  \t1 gwServes~g \neq \emptyset
\end{axdef}

The ids the operator may select on a gateway are the ones it is entitled to
serve --- the entitled set, not merely the catalogued set, so a
catalogued-but-unentitled id is refused at selection.

\begin{axdef}
  validOn : Gateway \fun \power ModelId
\where
  \forall g : Gateway @ validOn~g = gwServes~g
\end{axdef}

A model id's provider on a gateway, for the ids that gateway catalogs.

\begin{axdef}
  provOn : Gateway \cross ModelId \pfun Provider
\where
  \forall g : Gateway; m : ModelId @ \\
  \t1 ( ( g, m ) \in \dom provOn \iff m \in \dom (gwCatalog~g) ) \land \\
  \t1 ( m \in \dom (gwCatalog~g) \implies provOn~(g, m) = gwCatalog~g~m )
\end{axdef}

A gateway can carry these harnesses iff it exposes the OpenAI protocol.

\begin{axdef}
  usableGateways : \power Gateway
\where
  usableGateways = \{ g : Gateway | gwOpenai~g = yes \}
\end{axdef}

\section{Configuration state}

$sel$ is the selection; $selGw$ the gateway it was selected on; $routed$ whether
a gateway is active or the tools are on vendor defaults (direct); $activeGw$ the
gateway in force; $wire$ gives, per harness, per issued request kind, the id
that will travel. There is no override map and no auth variation: the OpenAI
harnesses route the id verbatim, so nothing needs bridging. $forced$ records a
\texttt{--force} selection so an off-entitlement id is representable.

\begin{schema}{Config}
  sel : ModelId \\
  selGw : Gateway \\
  routed : Bool \\
  activeGw : Gateway \\
  forced : Bool \\
  wire : Harness \pfun (ReqKind \pfun ModelId)
\where
  sel \in \dom (gwCatalog~selGw) \\
  forced = no \implies sel \in validOn~selGw
\end{schema}

\section{Initial state}

The never-configured state: direct, no managed wire.

\begin{schema}{Init}
  Config'
\where
  routed' = no \\
  forced' = no \\
  wire' = \{ h : \{ codex, pi, opencode \} @ h \mapsto \emptyset \}
\end{schema}

\section{Well-formedness}

$wire$ is defined for exactly the three harnesses and, for each, exactly the
request kinds it issues.

\begin{schema}{WellFormed}
  Config
\where
  \dom wire = \{ codex, pi, opencode \} \\
  ( \forall h : \dom wire @ \dom (wire~h) = issues~h )
\end{schema}

\section{The core guarantee}

\textbf{Routing correctness.} When routed, every request every harness issues
carries the selection verbatim --- no canonicalized stand-in, no unset-fallback
built-in. There is no conditional: main and side alike are the selection,
because no harness here rewrites the id.

\begin{schema}{RoutingCorrect}
  WellFormed
\where
  routed = yes \implies \\
  \t1 ( \forall h : \dom wire @ \forall k : \dom (wire~h) @ \\
  \t2   (wire~h)~k = sel )
\end{schema}

\textbf{Provider fidelity.} The wire id equals the selection, so the provider
reached equals the selection's provider on the active gateway. Because $validOn$
is the entitled set, that provider is one the account may actually reach ---
excluding the ``catalogued but \texttt{AccessDenied}'' failure by construction.

\begin{schema}{ProviderFidelity}
  RoutingCorrect
\where
  routed = yes \implies \\
  \t1 ( \forall h : \dom wire @ \forall k : \dom (wire~h) @ \\
  \t2   provOn~(activeGw, (wire~h)~k) = provOn~(activeGw, sel) )
\end{schema}

\textbf{The forbidden failure.} A side-request falling back to a built-in whose
provider is not the selection's. Named as the predicate the configuration makes
unreachable while routed.

\begin{schema}{SilentMisroute}
  WellFormed
\where
  routed = yes \\
  ( \exists h : \dom wire @ \exists k : \dom (wire~h) @ \\
  \t1 provOn~(activeGw, (wire~h)~k) \neq provOn~(activeGw, sel) )
\end{schema}

\section{Shared effect: applying a gateway configuration}

Route through gateway $g$ with selection $m$, rebuilding $wire$ so every issued
request kind of every harness maps to $m$.

\begin{schema}{ApplyGateway}
  \Delta Config \\
  g? : Gateway \\
  m? : ModelId \\
  fwd? : Bool
\where
  sel' = m? \\
  selGw' = g? \\
  activeGw' = g? \\
  routed' = yes \\
  forced' = fwd? \\
  \dom wire' = \{ codex, pi, opencode \} \\
  ( \forall h : \dom wire' @ wire'~h = \{ k : issues~h @ k \mapsto m? \} )
\end{schema}

\section{Operation: \texttt{use}}

Selecting a model. Input is a gateway, a model id, and a $force$ flag.
Partitions: (a) valid --- id entitled and gateway usable; (b) forced ---
$force? = yes$ records an off-entitlement id that is \emph{still catalogued}
(force bypasses the entitlement check, not the catalog: an id the gateway cannot
route at all is never selectable, which keeps every reachable selection's
provider well-defined), with the gateway usable; (c) refused-unknown --- id not
catalogued and not forced; (d) refused-gateway --- gateway does not expose the
OpenAI protocol.

\begin{schema}{UseOk}
  ApplyGateway \\
  force? : Bool
\where
  m? \in validOn~g? \\
  g? \in usableGateways \\
  fwd? = no \\
  force? = no
\end{schema}

\begin{schema}{UseForced}
  ApplyGateway \\
  force? : Bool
\where
  force? = yes \\
  fwd? = yes \\
  m? \in \dom (gwCatalog~g?) \\
  m? \notin gwServes~g? \\
  g? \in usableGateways
\end{schema}

\begin{schema}{UseRefusedUnknown}
  \Xi Config \\
  g? : Gateway \\
  m? : ModelId \\
  force? : Bool \\
  err! : \num
\where
  m? \notin \dom (gwCatalog~g?) \\
  force? = no \\
  err! = 1
\end{schema}

\begin{schema}{UseRefusedGateway}
  \Xi Config \\
  g? : Gateway \\
  m? : ModelId \\
  force? : Bool \\
  err! : \num
\where
  g? \notin usableGateways \\
  err! = 2
\end{schema}

\begin{zed}
  Use \defs (UseOk \lor UseForced) \lor \\
  \t1 (UseRefusedUnknown \lor UseRefusedGateway)
\end{zed}

\section{Operations: gateway-on, \texttt{direct}, \texttt{toggle}}

$GatewayOn$ re-applies the current selection on its gateway. $Direct$ withdraws
the managed wire, retaining $sel$/$selGw$ for a later re-apply. \texttt{toggle}
flips the route.

\begin{schema}{GatewayOn}
  \Delta Config
\where
  selGw \in usableGateways \\
  sel' = sel \\
  selGw' = selGw \\
  activeGw' = selGw \\
  routed' = yes \\
  forced' = forced \\
  \dom wire' = \{ codex, pi, opencode \} \\
  ( \forall h : \dom wire' @ wire'~h = \{ k : issues~h @ k \mapsto sel \} )
\end{schema}

\begin{schema}{Direct}
  \Delta Config
\where
  routed' = no \\
  sel' = sel \\
  selGw' = selGw \\
  activeGw' = activeGw \\
  forced' = forced \\
  wire' = \{ h : \{ codex, pi, opencode \} @ h \mapsto \emptyset \}
\end{schema}

\begin{schema}{Toggle}
  \Delta Config
\where
  sel' = sel \\
  selGw' = selGw \\
  forced' = forced \\
  ( routed = no \implies \\
  \t1 ( routed' = yes \land activeGw' = selGw \land \\
  \t2   selGw \in usableGateways \land \\
  \t2   \dom wire' = \{ codex, pi, opencode \} \land \\
  \t2   ( \forall h : \dom wire' @ wire'~h = \{ k : issues~h @ k \mapsto sel \} ) ) ) \\
  ( routed = yes \implies \\
  \t1 ( routed' = no \land activeGw' = activeGw \land \\
  \t2   wire' = \{ h : \{ codex, pi, opencode \} @ h \mapsto \emptyset \} ) )
\end{schema}

\section{Read-only operations: \texttt{status}, \texttt{list}, \texttt{verify}}

$Status$ reports the selection, gateway and route. $List$ reports the entitled
set. $Verify$ is the whole-fleet probe: $ok! = 1$ exactly when, routed and
well-formed (so the side-request slots exist), every harness's every request
kind reaches the selection's provider.

\begin{schema}{Status}
  \Xi Config \\
  repSel! : ModelId \\
  repGw! : Gateway \\
  repRouted! : Bool
\where
  repSel! = sel \\
  repGw! = activeGw \\
  repRouted! = routed
\end{schema}

\begin{schema}{List}
  \Xi Config \\
  cat! : \power ModelId
\where
  cat! = validOn~activeGw
\end{schema}

\begin{schema}{VerifyPass}
  \Xi Config \\
  WellFormed \\
  ok! : \num
\where
  routed = yes \\
  ( \forall h : \dom wire @ \forall k : \dom (wire~h) @ \\
  \t1 provOn~(activeGw, (wire~h)~k) = provOn~(activeGw, sel) ) \\
  ok! = 1
\end{schema}

\begin{schema}{VerifyFail}
  \Xi Config \\
  ok! : \num
\where
  ( routed = no \lor \\
  \t1 ( \lnot ( \dom wire = \{ codex, pi, opencode \} \land \\
  \t2     ( \forall h : \dom wire @ \dom (wire~h) = issues~h ) ) ) \lor \\
  \t1 ( \exists h : \dom wire @ \exists k : \dom (wire~h) @ \\
  \t2   provOn~(activeGw, (wire~h)~k) \neq provOn~(activeGw, sel) ) ) \\
  ok! = 0
\end{schema}

\begin{zed}
  Verify \defs VerifyPass \lor VerifyFail
\end{zed}

\section{The system invariant}

Whenever routed, the configuration is well-formed and every request carries the
selection and reaches its provider. Stated as a machine-checkable theorem so the
checker confirms it, not prose.

\begin{schema}{SystemOK}
  WellFormed
\where
  routed = yes \implies \\
  \t1 ( ( \forall h : \dom wire @ \forall k : \dom (wire~h) @ \\
  \t2     (wire~h)~k = sel ) \land \\
  \t1   ( \forall h : \dom wire @ \forall k : \dom (wire~h) @ \\
  \t2     provOn~(activeGw, (wire~h)~k) = provOn~(activeGw, sel) ) )
\end{schema}

\begin{theorem}{ApplyGivesSystemOK}
  \forall Config'; g? : Gateway; m? : ModelId; fwd? : Bool @ \\
  \t1 ( \exists Config @ ApplyGateway ) \implies \\
  \t2 ( \theta Config' \in \{ SystemOK @ \theta Config \} )
\end{theorem}

\section{Per-harness on-disk configuration}

A model slot holds a concrete id or $inherit$ (unset --- the harness falls back
to its parent / main model, which on the wire is the selection). A slot
\emph{resolves to the selection} when it is $inherit$ or $fixed~sel$. With no
canonicalization and no tier aliasing, this single predicate covers every slot
of every harness here.

\begin{zed}
  Slot ::= inherit | fixed \ldata ModelId \rdata
\end{zed}

\begin{axdef}
  resolves : ModelId \fun \power Slot
\where
  \forall m : ModelId @ \\
  \t1 resolves~m = \{ inherit \} \cup \{ s : Slot | s = fixed~m \}
\end{axdef}

\subsection{codex}

\texttt{\~{}/.codex/config.toml}: $cxMain$ (\texttt{model}) and the sub-agent
slot $cxSubagent$ (\texttt{[agents] agent\_default\_subagent\_model}). An omitted
sub-agent model inherits $cxMain$, so $inherit$ is the safe default.

\begin{schema}{CodexConfig}
  cxMain : Slot \\
  cxSubagent : Slot
\end{schema}

\begin{schema}{CodexOK}
  CodexConfig \\
  sel : ModelId
\where
  cxMain \in resolves~sel \\
  cxSubagent \in resolves~sel
\end{schema}

\subsection{pi}

\texttt{\~{}/.pi/agent/settings.json} ($piMain$, \texttt{defaultModel}); a
sub-agent Markdown file's omitted frontmatter \texttt{model:} inherits the
session model.

\begin{schema}{PiConfig}
  piMain : Slot \\
  piSubagent : Slot
\end{schema}

\begin{schema}{PiOK}
  PiConfig \\
  sel : ModelId
\where
  piMain \in resolves~sel \\
  piSubagent \in resolves~sel
\end{schema}

\subsection{opencode}

\texttt{\~{}/.config/opencode/opencode.json}: $ocMain$ (\texttt{model}),
$ocSmall$ (\texttt{small\_model}, pinned to the selection), and per-agent slots
$ocAgent$ (\texttt{agent.<name>.model}, inherit-safe).

\begin{schema}{OpencodeConfig}
  ocMain : Slot \\
  ocSmall : Slot \\
  ocAgent : Slot
\end{schema}

\begin{schema}{OpencodeOK}
  OpencodeConfig \\
  sel : ModelId
\where
  ocMain \in resolves~sel \\
  ocSmall \in resolves~sel \\
  ocAgent \in resolves~sel
\end{schema}

\subsection{The composed fleet invariant}

When routed, every harness --- main and every side slot --- resolves to the
selection. The file-level restatement of $SystemOK$.

\begin{schema}{Fleet}
  sel : ModelId \\
  routed : Bool \\
  CodexConfig \\
  PiConfig \\
  OpencodeConfig
\end{schema}

\begin{schema}{FleetOK}
  Fleet
\where
  routed = yes \implies ( CodexOK \land PiOK \land OpencodeOK )
\end{schema}

\section{Why this tool is simple}

Everything that made the combined tool fragile lived on the Claude Code path and
is absent here: no canonicalization ($wireOf$ would be identity, so it is not
even defined), no override map, no tier-alias fan-out, no $[1m]$ marker, no
subscription-credential mode, one protocol. What remains is: write the verbatim
selection into three config files whose main and side slots all take the same
id. The invariant is a single equality ($(wire~h)~k = sel$ everywhere), which is
why the whole tool model-checks to one uniform guarantee with no per-request
case analysis.

\end{document}
